监控面板上，某个指标的日均值三个月纹丝不动。团队据此判断系统稳定。

然后有人调出了逐分钟的曲线：它在两个水平之间反复横跳，每次切换持续几个小时。日均值确实没变——它是两个水平的加权平均，而权重恰好稳定。可这个平均值从来没有出现过，也没有任何一分钟的系统状态与它对应。

概率论到这里就够不着了。它能描述某一个时刻的分布，却不描述\emph{同一次实验里，相邻时刻怎样彼此牵连}。给随机变量加上时间下标，多出来的不是一个记号，是一整套新问题：一条轨迹会不会回到出发点、会不会永远漂走、需要记住多少历史才能预测下一步、以及最要紧的那个——观察一条足够长的轨迹，能不能替代观察所有可能的轨迹。

最后这个问题值得单独摆出来，因为几乎所有从时间序列做估计的做法都押在它上面。

\begin{center}
\begin{tikzpicture}[x=1cm,y=1cm,font=\small,>=stealth]
  \draw[->,black!55] (0,0) -- (6.6,0) node[right,font=\footnotesize] {\(t\)};
  \draw[->,black!55] (0,-0.2) -- (0,2.9);
  \draw[blue!65,thick] plot[smooth] coordinates
    {(0.2,2.25)(0.7,2.05)(1.2,2.35)(1.7,2.1)(2.2,2.3)(2.7,2.15)(3.2,2.4)(3.7,2.1)(4.2,2.3)(4.7,2.2)(5.2,2.35)(5.7,2.1)(6.2,2.25)};
  \draw[orange!80!black,thick] plot[smooth] coordinates
    {(0.2,0.75)(0.7,0.95)(1.2,0.65)(1.7,0.9)(2.2,0.7)(2.7,0.85)(3.2,0.6)(3.7,0.9)(4.2,0.7)(4.7,0.8)(5.2,0.65)(5.7,0.9)(6.2,0.75)};
  \draw[black!45,dashed] (0,1.5) -- (6.4,1.5);
  \node[right,font=\footnotesize,black!60] at (6.45,1.5) {集合平均};
  \node[right,font=\footnotesize,blue!65] at (6.45,2.2) {轨道 A};
  \node[right,font=\footnotesize,orange!80!black] at (6.45,0.78) {轨道 B};
  \draw[black!40] (2.9,-0.35) -- (2.9,2.75);
  \node[below,font=\footnotesize,black!60] at (2.9,-0.4) {某个时刻的截面};
  \node[align=center,font=\footnotesize,black!70] at (9.6,1.5)
    {沿轨道求时间平均：\(2.2\) 或 \(0.8\)\\[2pt]沿截面求集合平均：\(1.5\)\\[2pt]\textbf{没有一条轨道到过 \(1.5\)}};
\end{tikzpicture}

\emph{每个时刻的截面分布都不随时间改变，所以这个过程是平稳的；但时间平均与集合平均不相等，所以它不遍历。用一条长轨迹估计期望的做法，在这里给出的答案取决于你抽到了哪一条。}
\end{center}

\subsection{五条贯穿始终的判断}

\textbf{先分清是在沿轨道看，还是沿截面看。}固定一次实验看一整条历史，与固定一个时刻看所有可能世界，是两个方向。初学的困惑大多来自把它们混为一谈——说``这个过程收敛了''，必须先说清是哪一种。只跑一次训练就下结论，等于用一条轨道回答一个截面问题。

\textbf{平稳保证极限存在，遍历才保证极限是常数。}上面那张图就是反例：平稳而不遍历时，时间平均是个随机量，取决于你落在哪条轨道上。MCMC 的合法性依赖的正是遍历性而不是样本独立性；采样器没跨过模态时，你拿到的就是图里那两条互不相通的轨道之一。

\textbf{Markov 性质是关于状态设计的断言，不是关于系统的。}同一个系统，状态取得太粗就不是 Markov 的，把必要的历史并进状态就又是了。所以遇到``这个过程不满足 Markov 性''，该动的往往是状态定义而不是模型族。强化学习里状态是否充分、序列模型的上下文窗口截断多少，问的都是这件事。

\textbf{长期行为由谱决定，收敛速度由谱间隙决定。}转移矩阵的最大特征值是 1，对应平稳分布；第二大特征值的模决定每步误差衰减多少，也就决定了 burn-in 要多久。这与线性代数里矩阵幂的长期行为、递推关系的特征根、线性系统的稳定性是同一件事，只是换了应用场合。连续时间里矩阵幂换成矩阵指数 \(e^{tQ}\)，判据也随之从``模小于 1''变成``实部为负''。

\textbf{到达率低于服务率只是稳定的门槛，不是良好的保证。}利用率趋近 1 时队长按 \(\rho/(1-\rho)\) 爆炸——余量从 \(20\%\) 减到 \(10\%\)，等待时间直接翻倍。容量规划按平均负载做，就会撞上这条曲线的陡峭段。

\subsection{七章怎么安排}

\hyperref[ch:046]{``随机过程基础：从一个随机数到一整条随机轨迹''}是入口，轨道与截面的区分、平稳与遍历的落差都在这里，值得慢读。

\hyperref[ch:048]{``离散时间 Markov 链：状态足够好，历史就可以被压缩''}是全部分用得最多的一章，也是后面强化学习、MCMC、序列模型的共同前置。它和\hyperref[ch:047]{``计数过程：事件何时到来、到来了多少''}构成两条基本模型线——前者管状态怎么转移，后者管事件何时发生。

\hyperref[ch:050]{``连续时间 Markov 链与排队：不仅要知道跳到哪，还要知道何时跳''}把离散时间的结论搬到连续时间，其中 Little 定律几乎不依赖任何分布假设，是实践中性价比最高的一条结果。

\hyperref[ch:051]{``扩散与随机微分方程：路径不可微时怎样写微分方程''}在本部分里最接近当代生成模型：前向加噪、反向去噪、得分函数这条线就是扩散模型的数学地基。想直接理解扩散模型的读者可以从这一章切入，需要的前置只有随机游走与布朗运动。

\hyperref[ch:049]{``随机过程的四种典型用法''}与\hyperref[ch:052]{``ARMA 与 ARIMA：时间序列的经典模型族''}偏应用，可按需取用。前者里的可选停时定理值得一读——它说明没有任何停止策略能把公平游戏变成有利可图，而这正解释了``看到显著就停''为什么会让实验失效。

读这一部分时保持一个追问：手上这个结论是关于\emph{一条轨迹}的，还是关于\emph{所有轨迹的分布}的？两者需要的证据不同，能支撑的判断也不同。混淆它们不会报错，只会让人对一个从未出现过的平均值充满信心。
